\documentclass{article}
\usepackage[T1]{fontenc}
\usepackage{lmodern}
\usepackage{graphicx}
\usepackage{amsmath,amssymb}
\usepackage[a4paper,margin=2cm]{geometry}
\usepackage[absolute,overlay]{textpos}
\setlength{\TPHorizModule}{1mm}
\setlength{\TPVertModule}{1mm}
\usepackage[indonesian]{babel}
\usepackage{hyperref}
\setlength{\emergencystretch}{2em}
% unit: Halaman 89

\begin{document}

% === Halaman 89 (llm) ===

\textbf{6. \textit{Keystream} dapat dihasilkan secara paralel}

Salah satu karakteristik penting Salsa20 adalah setiap blok \textit{keystream} bergantung pada \textit{key}, \textit{nonce}, dan \textit{counter}.

Secara konseptual:

\begin{align*}
\text{Key + Nonce + Counter 0} & \rightarrow \text{Keystream 0} \\
\text{Key + Nonce + Counter 1} & \rightarrow \text{Keystream 1} \\
\text{Key + Nonce + Counter 2} & \rightarrow \text{Keystream 2} \\
\text{Key + Nonce + Counter 3} & \rightarrow \text{Keystream 3}
\end{align*}

Blok-blok tersebut pada prinsipnya dapat dihitung secara independen. Ini berbeda dari beberapa konstruksi \textit{stream cipher} yang memiliki ketergantungan berantai antar keystream.

\end{document}
